%%==========================================================================
\section{Introduction}
\label{sec:intro}
%%==========================================================================

The airline planning process is traditionally decomposed into a sequence of
sub-problems: timetable design, fleet assignment, crew rostering, tail
assignment, and recovery planning~\citep{barnhart2003applications,
gopalan1998mathematical}.  The \emph{tail assignment problem} (TAP)—assigning
a specific aircraft (identified by its tail number) to each flight leg—is the
penultimate planning step and is tightly coupled with aircraft maintenance
scheduling.  Poor tail assignment decisions directly inflate operational costs
through unnecessary ferry (repositioning) flights, avoidable leasing decisions,
suboptimal maintenance routing, and service disruptions.

\citet{khaled2018compact} proposed a compact 0--1 linear programming
formulation for the TAP, featuring only $O(n_f n_p)$ binary variables (where
$n_f$ is the number of flights and $n_p$ the fleet size).  This is a
significant reduction over the classic multi-commodity flow model of
\citet{levin1971scheduling}, which requires $O(n_f^2 n_p)$ variables.
Computational experiments reported in that paper show that instances with up
to 1\,500 flight legs and 40 aircraft can be solved to certified optimality
within one hour on a standard workstation.

However, most practical studies still rely on cyclic or short-horizon settings
(daily/weekly schedules) and often separate routing from maintenance decisions.
For modern airline operations, this is restrictive: schedules are frequently
adjusted for seasonality and event-driven demand, while aircraft must satisfy
multi-level maintenance obligations across the same horizon.  Treating routing
and maintenance independently shifts feasibility issues into the operational
window, where corrective actions are expensive and disruptive.

This thesis addresses that gap by revisiting and extending the compact TAP
framework toward an integrated, season-level routing-and-maintenance planning
approach.  The objective is to produce physically executable aircraft rotations
that satisfy maintenance requirements and improve planning certainty before
schedule execution.

\subsection*{Thesis contributions}

This thesis makes four contributions that verify, correct, and extend the
state of the art:

\begin{enumerate}
\item \textbf{Basic feasibility correction (C2--C4 / constraints (3)--(6)).}
  We show that the strict-time balance definition used in the compact basic
  model admits a simultaneous-departure corner case: one aircraft can be
  assigned to two flights departing at the same airport and same time without
  violating the balance inequalities.  We formalize this counter-example and
  enforce physical executability by adding explicit pairwise non-overlap
  constraints.

\item \textbf{Constraint~(13) correction.}  We show that the cumulative-flight-
  time maintenance constraint (Constraint~13 in \citet{khaled2018compact})
  contains a logical flaw: the big-$M$ linearisation term $M(2 - y_{jd} -
  y_{jd'})$ leaves the constraint vacuously satisfied whenever at least one of
  the endpoint days $d, d'$ hosts no check, so an accumulation window that is
  not bounded by checks at both ends is left unconstrained and may exceed the
  flight-hour limit.  We provide a corrected formulation using two split
  endpoint constraints, prove its validity, and show that the corrected
  formulation yields strictly tighter LP relaxations on high-density instances.

\item \textbf{Full $\{A, B, C, D\}$ maintenance hierarchy.}  The model of
  \citet{khaled2018compact} handles type-A checks (cumulative flight hours).
  We extend the formulation to the full four-level hierarchy mandated by
  aviation regulators (FAA/EASA): type-A and~B checks are flight-hour
  triggered; type-C and~D checks are calendar-day triggered.  A heavier check
  subsumes all lighter checks (hierarchical counter resets).  We further add
  a constraint accounting for pre-horizon accumulated maintenance hours/days per
  aircraft and integrate these requirements into long-horizon schedule planning.

\item \textbf{Heuristic benchmark.}  We implement and evaluate two
  constructive heuristics: (a) a greedy/insertion heuristic operating
  directly on flight lists, and (b) a Dijkstra-based heuristic on the
  space-time network representation (adapted from
  \citet{gronkvist2005tail}).  Both heuristics respect the full maintenance
  hierarchy.  We report optimality gaps against the MILP on the same instance
  sets and show the Dijkstra heuristic consistently outperforms an ACO
  alternative.
\end{enumerate}

\subsection*{Thesis organisation}

\Cref{ch:lit} reviews related work.  \Cref{ch:ground} restates the
basic TAP model and proves a C2--C4 corner-case correction.  \Cref{ch:maintenance}
presents the maintenance-extended model, the Constraint~(13) correction, and
the induced interpretation updates for C8--C14 under the corrected basic model.
\Cref{ch:hierarchy} describes the
full $\{A,B,C,D\}$ extension.  \Cref{ch:heuristics} analyses the heuristic
algorithms.  \Cref{ch:experiments} reports all computational results.
\Cref{ch:conclusion} concludes.
